% ============================================================
% UniConnecT — Project Proposal
% CSE 3422: Software Engineering Laboratory, Summer 2026
% United International University, Dhaka
% Style follows the UIU CSE (FYDP) LaTeX conventions:
% A4, 11pt, 1.5 line spacing, ~2.5cm margins, formal front matter.
% Compile with pdflatex (Overleaf-ready).
% ============================================================
\documentclass[11pt, a4paper]{article}

\usepackage[a4paper, margin=2.5cm]{geometry}
\usepackage{setspace}
\usepackage{graphicx}
\usepackage{booktabs}
\usepackage{enumitem}
\usepackage{array}
\usepackage{xcolor}
\usepackage{pgfgantt}
\usepackage{fancyhdr}
\usepackage[hidelinks]{hyperref}
\usepackage{titlesec}

\definecolor{uiuorange}{RGB}{246, 139, 31}
\definecolor{deepnavy}{RGB}{16, 24, 40}

\titleformat{\section}{\large\bfseries\color{deepnavy}}{\thesection}{0.8em}{}
\titleformat{\subsection}{\normalsize\bfseries\color{deepnavy}}{\thesubsection}{0.8em}{}

\pagestyle{fancy}
\fancyhf{}
\fancyhead[L]{\small UniConnecT --- Project Proposal}
\fancyhead[R]{\small CSE 3422, Summer 2026}
\fancyfoot[C]{\small\thepage}
\renewcommand{\headrulewidth}{0.4pt}

\onehalfspacing

\begin{document}

% ------------------------------------------------------------
% Title page
% ------------------------------------------------------------
\begin{titlepage}
    \centering
    % Place the UIU logo file (uiu_logo.png) next to this .tex file
    % \includegraphics[width=3.5cm]{uiu_logo.png}\\[0.6cm]
    {\Large \textbf{United International University}}\\[0.2cm]
    {\large Department of Computer Science and Engineering}\\[2cm]

    {\large \textbf{Project Proposal}}\\[0.4cm]
    {\huge \textbf{UniConnecT}}\\[0.3cm]
    {\large A Multi-Tenant Private Social Network and Campus Platform\\for Universities}\\[2cm]

    {\normalsize
    \begin{tabular}{ll}
        \textbf{Course}    & CSE 3422: Software Engineering Laboratory \\
        \textbf{Trimester} & Summer 2026 \\
        \textbf{Section}   & \texttt{<section>} \\
    \end{tabular}
    }\\[1.5cm]

    {\large \textbf{Submitted by --- Team Mavericks}}\\[0.3cm]
    {\normalsize
    \begin{tabular}{lll}
        \toprule
        \textbf{Name} & \textbf{Student ID} & \textbf{Role} \\
        \midrule
        Joydip Datta          & 011~233~0960        & Team lead, full-stack \\
        \texttt{<member 2>}   & \texttt{<id>}       & \texttt{<role>} \\
        \texttt{<member 3>}   & \texttt{<id>}       & \texttt{<role>} \\
        \texttt{<member 4>}   & \texttt{<id>}       & \texttt{<role>} \\
        \bottomrule
    \end{tabular}
    }\\[1.5cm]

    {\normalsize \textbf{Submitted to}}\\[0.2cm]
    {\normalsize \texttt{<Instructor name>}\\
    \texttt{<Designation>}, Department of CSE\\
    United International University}\\[1.2cm]

    {\normalsize Date of submission: \texttt{<date>}}
    \vfill
\end{titlepage}

% ------------------------------------------------------------
\section{Introduction}

University communities in Bangladesh coordinate their academic and social
lives across a fragmented set of general-purpose tools: WhatsApp groups for
class coordination, Facebook pages for club announcements, LinkedIn for
alumni networking, and email for official notices. Information is scattered,
unverified accounts mix with real ones, and nothing is aware of the
distinction between a student, a faculty member, and an alumnus.

\textbf{UniConnecT} is a private, institution-scoped social network built
specifically for universities. Every account is verified through an
invite-and-OTP flow against the university's own email domain, every feature
is role-aware (student, alumni, faculty, admin, and transport staff), and the
platform is architected as multi-tenant SaaS so any university can be
onboarded as an isolated tenant.

A functional first version of UniConnecT was designed, implemented, and
\emph{deployed to production} by our team in the Spring 2026 trimester for
the Web Programming course, where it grew substantially beyond the required
feature set (Section~\ref{sec:existing}). In Summer 2026, under the Software
Engineering Laboratory, we propose to take the project through a full
software-engineering lifecycle: formal requirements specification, design
documentation with UML, disciplined agile iterations, systematic testing and
quality assurance, and a set of new high-impact features that prepare the
platform for real adoption at UIU and beyond.

\section{Problem Statement}

\begin{enumerate}[leftmargin=1.5em, itemsep=2pt]
    \item \textbf{Fragmentation.} Campus information lives in dozens of
    unofficial WhatsApp/Facebook groups with no single source of truth;
    important notices are routinely missed.
    \item \textbf{No identity verification.} Public platforms cannot
    guarantee that a ``UIU student'' is actually a UIU student, enabling
    impersonation and spam.
    \item \textbf{No role awareness.} Faculty announcements, alumni job
    referrals, and student discussions all compete in the same unstructured
    channels.
    \item \textbf{Weak alumni linkage.} Universities lose contact with
    graduates, and students lose access to mentorship and referral networks.
    \item \textbf{Scattered campus services.} Shuttle schedules, lost-and-found,
    course materials, and event registration have no common digital home.
\end{enumerate}

\section{Objectives}

\begin{enumerate}[leftmargin=1.5em, itemsep=2pt]
    \item Apply the full software engineering process --- SRS, UML-based
    design, agile sprint planning, code review, CI/CD, and multi-level
    testing --- to a real, deployed product.
    \item Polish the existing platform to production quality: accessibility
    (WCAG 2.1 AA), performance on mid-tier mobile devices, security
    hardening, and end-to-end test coverage.
    \item Deliver a set of new differentiating features (Section~\ref{sec:proposed}),
    headlined by an AI campus assistant and a self-serve university
    onboarding flow that completes the multi-tenant SaaS vision.
    \item Publish the platform for broader use --- a public landing/marketing
    site, demo tenant for evaluation, and documentation that allows a new
    university to be onboarded without developer involvement.
\end{enumerate}

% ------------------------------------------------------------
\section{Work Completed in Spring 2026 (Existing System)}
\label{sec:existing}

The Web Programming course required a core social-network feature set, all
of which were delivered: user registration and profiles, status posts,
comments, connection (friend) requests, activity notifications, a responsive
newsfeed, a like/reaction system, profile picture upload, user search, and a
database-driven activity log.

Beyond that requirement, the team continued development and shipped the
following major subsystems, all live in production today:

\begin{description}[leftmargin=1.5em, itemsep=2pt]
    \item[Platform \& security] Multi-tenant architecture with strict
    per-university data isolation; five roles with least-privilege access;
    invite-only registration with email OTP; JWT access tokens with rotating
    refresh sessions; per-university audit logging.
    \item[Real-time layer] Socket.io messaging (direct, group, and mentorship
    conversations), online-presence tracking with privacy tiers, and live
    web-push notifications --- scaled across instances via Redis pub/sub.
    \item[Career \& community] Job board, events with registration, groups
    (private membership, moderator roles, shared resources, study sessions),
    and a LinkedIn-style profile suite (experience, education, featured work,
    profile analytics and viewers, vanity URLs, resume export).
    \item[Mentorship economy] Alumni--student mentorship with request
    lifecycle automation, mentor capacity limits, session logging, and a
    points-based reward system redeemable as gift cards.
    \item[Academic tools] A lightweight LMS (course outlines, modules,
    assignments, gradebook), a quiz engine with answer review, and
    gamified learning streaks and badges.
    \item[Campus services] Lost-and-found, shuttle schedules with
    \emph{live GPS bus tracking} broadcast by a dedicated driver role, and
    automated import of official university notices from the university
    website (WordPress API with browser-automation fallback).
    \item[Discovery \& ranking] Postgres full-text search across people,
    posts, jobs, events, and groups; a trending ``Top'' feed with a
    Hacker-News-style hot score recomputed by background workers; hashtag
    explore pages; drafts and scheduled publishing.
    \item[Administration] Admin dashboard with university statistics, user
    and role management, bulk invitations, content-report moderation, and
    content-sync review queues.
    \item[Infrastructure] Monorepo (React + Express + shared type-safe
    contracts), eight background job queues, presigned direct-to-storage
    file uploads, and a production deployment on Azure (API), Vercel (web),
    and Neon (Postgres) with CI/CD from GitHub.
\end{description}

% ------------------------------------------------------------
\section{Proposed Work for Summer 2026}
\label{sec:proposed}

The Summer 2026 scope has three tracks: \emph{engineering rigor},
\emph{polish for publication}, and \emph{new features}.

\subsection{Track A --- Software engineering rigor (course deliverables)}
\begin{itemize}[leftmargin=1.5em, itemsep=2pt]
    \item Software Requirements Specification (IEEE 830 style) covering the
    full platform.
    \item UML design set: use-case, class, sequence, activity, and ER
    diagrams; C4 architecture diagrams.
    \item Agile process: 2-week sprints with sprint planning, backlog
    grooming, burndown tracking, and retrospectives (GitHub Projects).
    \item Quality gates: mandatory code review, CI pipeline running lint,
    type checks, unit and integration tests on every pull request.
    \item Test plan and traceability matrix; end-to-end test suite
    (Playwright) for the ten highest-value user journeys; target
    $\geq 70\%$ coverage on service-layer code.
\end{itemize}

\subsection{Track B --- Polish and publication}
\begin{itemize}[leftmargin=1.5em, itemsep=2pt]
    \item \textbf{Accessibility audit} to WCAG 2.1 AA (contrast, keyboard
    navigation, screen-reader roles, reduced motion).
    \item \textbf{Performance hardening} for mid-tier Android devices:
    Lighthouse $\geq 90$, bundle-size budget, image optimization, API
    response caching.
    \item \textbf{Progressive Web App}: installable app with offline shell
    and cached feed, completing the existing push-notification support.
    \item \textbf{Security review}: rate limiting on sensitive endpoints,
    dependency scanning, OWASP Top-10 checklist, penetration-test style
    review of the multi-tenant isolation boundary.
    \item \textbf{Observability}: structured error tracking, uptime and
    performance dashboards, alerting on production incidents.
    \item \textbf{Localization groundwork}: Bangla language support for the
    highest-traffic surfaces.
    \item \textbf{Public launch kit}: marketing landing page, demo tenant
    with seeded data for evaluators, and administrator documentation.
\end{itemize}

\subsection{Track C --- New features (differentiators)}
\begin{enumerate}[leftmargin=1.5em, itemsep=2pt]
    \item \textbf{AI campus assistant.} A retrieval-augmented chatbot that
    answers questions over university notices, events, shuttle schedules,
    and course materials (``When does the Summer registration close?'',
    ``Summarize this week's CSE notices''), plus AI-assisted post
    summarization and semantic search. This is the headline feature.
    \item \textbf{Self-serve university onboarding.} A guided flow where a
    university administrator registers a domain, verifies ownership,
    configures branding, and bulk-invites members --- turning the
    multi-tenant architecture into a true SaaS product.
    \item \textbf{Smart recommendations.} ``People you may know,'' job and
    event recommendations driven by profile, department, and connection-graph
    signals.
    \item \textbf{Event check-in with QR codes.} Organizers scan attendee QR
    passes; attendance feeds into event analytics.
    \item \textbf{Resume builder.} One-click, professionally formatted PDF
    resume generated from the structured profile (experience, education,
    skills, featured work).
    \item \textbf{Admin analytics dashboard.} Engagement trends, retention
    cohorts, and content-health metrics for university administrators.
\end{enumerate}

% ------------------------------------------------------------
\section{Target Users and Stakeholders}

\begin{center}
\begin{tabular}{p{3.2cm}p{10.5cm}}
    \toprule
    \textbf{Stakeholder} & \textbf{Value delivered} \\
    \midrule
    Students & Verified campus feed, groups and study sessions, jobs,
    mentorship, shuttle tracking, course tools. \\
    Alumni & Lifelong institutional identity, mentorship rewards, referral
    and hiring channel. \\
    Faculty & Role-aware announcements, course material distribution,
    moderated discussion. \\
    University admin & Verified community, moderation and audit tooling,
    engagement analytics, notice distribution. \\
    Transport staff & Single-purpose GPS broadcast app for live shuttle
    tracking. \\
    \bottomrule
\end{tabular}
\end{center}

\section{Methodology}

We follow \textbf{Scrum} with six two-week sprints across the trimester.
Each sprint produces a potentially shippable increment deployed to a staging
environment; production releases occur after sprint review. Requirements are
maintained as a prioritized product backlog derived from the SRS; every user
story carries acceptance criteria that map to automated tests. GitHub flow
(feature branches, pull requests, mandatory review, CI gates) governs all
changes, and the definition of done includes updated documentation and
passing end-to-end tests.

\section{Tools and Technologies}

\begin{center}
\begin{tabular}{p{4cm}p{9.7cm}}
    \toprule
    \textbf{Layer} & \textbf{Technology} \\
    \midrule
    Frontend & React 18, TypeScript, Vite, Tailwind CSS, TanStack Query,
    Zustand, Socket.io client \\
    Backend & Node.js, Express, TypeScript, Socket.io, Zod, Knex, Bull
    queues \\
    Data \& infra & PostgreSQL (Neon), Redis, S3-compatible object storage,
    Azure App Service, Vercel \\
    AI features & LLM API with retrieval-augmented generation over
    tenant-scoped content, pgvector embeddings \\
    Engineering & GitHub Actions CI/CD, Playwright, Vitest, ESLint,
    Lighthouse CI, draw.io/PlantUML for UML \\
    Process & GitHub Projects (Scrum board), LaTeX (documentation) \\
    \bottomrule
\end{tabular}
\end{center}

\section{Feasibility}

\textbf{Technical.} The riskiest foundations --- multi-tenancy, real-time
infrastructure, background processing, and production deployment --- are
already built and battle-tested in production, so Summer work concentrates
on well-understood additions. The AI assistant uses managed LLM APIs rather
than custom model training, keeping complexity and cost bounded.
\textbf{Economic.} All infrastructure runs on free or student-tier cloud
plans; the only recurring cost is metered LLM usage, which is capped
per-tenant. \textbf{Operational.} The team has three trimesters of shared
history on this codebase and an established CI/CD pipeline, minimizing
process risk. \textbf{Schedule.} The plan reserves the final sprint
exclusively for stabilization, documentation, and presentation preparation.

\section{Work Plan and Timeline}

\begin{center}
\begin{ganttchart}[
    hgrid, vgrid,
    time slot format=isodate,
    time slot unit=day,
    x unit=0.105cm,
    y unit title=0.6cm,
    y unit chart=0.55cm,
    title label font=\scriptsize,
    bar label font=\scriptsize,
    group label font=\scriptsize\bfseries,
    milestone label font=\scriptsize\itshape,
    bar/.append style={fill=uiuorange!70},
    group/.append style={fill=deepnavy!80},
]{2026-06-01}{2026-08-31}
    \gantttitlecalendar{month=name} \\
    \ganttgroup{Sprint 1--2: SRS, UML, test infra}{2026-06-01}{2026-06-27} \\
    \ganttbar{SRS + UML design set}{2026-06-01}{2026-06-20} \\
    \ganttbar{E2E test suite + CI gates}{2026-06-08}{2026-06-27} \\
    \ganttgroup{Sprint 3--4: New features}{2026-06-28}{2026-07-25} \\
    \ganttbar{AI campus assistant}{2026-06-28}{2026-07-25} \\
    \ganttbar{Self-serve onboarding, QR check-in}{2026-06-28}{2026-07-18} \\
    \ganttbar{Recommendations, resume builder}{2026-07-05}{2026-07-25} \\
    \ganttgroup{Sprint 5: Polish \& publish}{2026-07-26}{2026-08-15} \\
    \ganttbar{A11y, performance, security, PWA}{2026-07-26}{2026-08-15} \\
    \ganttbar{Analytics dashboard, launch kit}{2026-07-26}{2026-08-15} \\
    \ganttgroup{Sprint 6: Stabilize \& present}{2026-08-16}{2026-08-31} \\
    \ganttbar{Testing, docs, final report}{2026-08-16}{2026-08-28} \\
    \ganttmilestone{Final presentation}{2026-08-30}
\end{ganttchart}
\end{center}

\noindent\emph{Adjust sprint dates to the official Summer 2026 academic calendar.}

\section{Team and Work Distribution}

\begin{center}
\begin{tabular}{p{4cm}p{9.7cm}}
    \toprule
    \textbf{Member} & \textbf{Responsibilities} \\
    \midrule
    Joydip Datta & Architecture, AI assistant, CI/CD, deployment, code review \\
    \texttt{<member 2>} & \texttt{<e.g., SRS + UML documentation, onboarding flow>} \\
    \texttt{<member 3>} & \texttt{<e.g., testing lead, QR check-in, analytics>} \\
    \texttt{<member 4>} & \texttt{<e.g., accessibility, PWA, resume builder>} \\
    \bottomrule
\end{tabular}
\end{center}

\section{Expected Outcome}

By the end of Summer 2026 we will deliver: (1) a complete SE documentation
set (SRS, UML design, test plan, traceability matrix); (2) a production
platform hardened to WCAG 2.1 AA and Lighthouse $\geq 90$ with automated
end-to-end coverage of critical journeys; (3) six new features headlined by
a tenant-scoped AI campus assistant and self-serve university onboarding;
and (4) a public launch kit enabling any university to evaluate and adopt
UniConnecT. The result is not a classroom prototype but a publishable,
multi-tenant SaaS product engineered with verifiable process discipline.

\section{Conclusion}

UniConnecT began as a course project and has already outgrown that framing:
it is a deployed, multi-tenant campus platform in production use. The
Software Engineering Laboratory is the natural next step --- applying formal
requirements, design, process, and quality engineering to carry the product
from ``working'' to ``publishable.'' We believe the combination of a real
deployed system, disciplined engineering artifacts, and ambitious new
features positions this project to be a standout submission for Summer 2026.

\begin{thebibliography}{9}
\bibitem{sommerville} I.~Sommerville, \emph{Software Engineering}, 10th ed., Pearson, 2015.
\bibitem{ieee830} IEEE Std 830-1998, \emph{IEEE Recommended Practice for Software Requirements Specifications}.
\bibitem{scrum} K.~Schwaber and J.~Sutherland, \emph{The Scrum Guide}, 2020.
\bibitem{wcag} W3C, \emph{Web Content Accessibility Guidelines (WCAG) 2.1}, 2018.
\bibitem{repo} Team Mavericks, \emph{UniConnecT source repository}, GitHub, 2026.
\end{thebibliography}

\end{document}
